\section{杂项}
\subsection{动态大小bitset}
\begin{lstlisting}[caption={}]
template<int N>
void solve(int n){
    while(N <= n){
        solve<std::min(N << 1,max_N)>(n);
        return 0;
    }
    std::bitset<N> st;
}
\end{lstlisting}

\subsection{Gosper's Hack}
生成$n$元集合所有$k$元子集的算法,只会遍历$C_n^k$种情况。
\begin{lstlisting}[caption={}]
void GospersHack(int k, int n){
    int cur = (1 << k) - 1;
    int limit = (1 << n);
    while (cur < limit && cur){
        // do something
        int lb = cur & -cur;
        int r = cur + lb;
        cur = ((r ^ cur) >> __builtin_ctz(lb) + 2) | r;
        // 或：cur = (((r ^ cur) >> 2) / lb) | r;
    }
}
\end{lstlisting}

\subsection{绝对众数}
\subsubsection{摩尔投票}
选出$1$个候选人，要求这个人的得票都超过总票数的一半 
\begin{lstlisting}[caption={}]
int m = 0,cnt = 0;
for(int i = 0;i < n;i++){
    if(cnt == 0) m = nums[i];
    if(m == nums[i]) cnt++;
    else cnt--;
}
// 最后需要验证答案是否符合要求
\end{lstlisting}

\subsubsection{拓展摩尔投票}
选出$N$个候选人，要求这个人的得票都超过总票数的$\frac{1}{N + 1}$
\begin{lstlisting}[caption={}]
int m[N],cnt[N];
for(auto e : nums){
    int i = find(m, m + N, e) - m;
    if(i != N){ // 如果当前票投给了候选人之一
        cnt[i]++;
        continue;
    }
    int j = find(cnt, cnt + N, 0) - cnt;
    if(j != N){ // 如果当前存在一个位置"虚位以待"
        m[j] = e;
        cnt[j] = 1;
        continue;
    }
    for(auto &c : cnt) c--;
}
// 最后需要验证答案是否符合要求
\end{lstlisting}

\subsubsection{区间绝对众数(线段树上摩尔投票)}
当绝对众数存在时，可以任意交换摩尔投票的计票顺序，而不改变选出的候选人。
在合并左右两个区间时，相当于是先对左右两边分别进行摩尔投票，再对两边抵消后的结果进行摩尔投票。
\begin{lstlisting}[caption={}]
struct Node{
    int m,cnt;
    Node operator+(const Node &o) const{
        if (m == o.m) return {m, cnt + o.cnt};
        else if (cnt > o.cnt) return {m, cnt - o.cnt};
        return {o.m, o.cnt - cnt};
    }
};
// 最后需要验证答案是否符合要求
// vector<int> pos[MAXV];
// pos[m]表示权值m的下标数组，需要预处理
// 如果权值值域过大，需要改用map或哈希表
auto itr = upper_bound(pos[m].begin(), pos[m].end(), right);
auto itl = lower_bound(pos[m].begin(), pos[m].end(), left);
if (itr - itl > (right - left + 1) / 3)
    return m;
else
    return -1;
\end{lstlisting}

\subsection{模拟退火}
\begin{lstlisting}[caption={}]
using point = std::array<double, 2>;
using pointW = std::array<double, 3>;
const double T0 = 10000.0; //初始温度
const double q = 0.98; //退火系数
const int L = 100; //每次迭代次数
std::mt19937_64 rng{std::chrono::steady_clock::now().time_since_epoch().count()};
std::uniform_real_distribution<double> rnd(-1.0, 1.0);
point Point{}, bestPoint{};
double best = 1e18;
std::vector<pointW> points;
double dist(const point &a, const pointW &b) {
    return sqrtl((a[0] - b[0]) * (a[0] - b[0]) + (a[1] - b[1]) * (a[1] - b[1])) * b[2];
}
double calc(const point &p) {
    double res = 0;
    for (int i = 0; i < (int)points.size(); i++)
        res = std::max(res, dist(p, points[i]));
    return res;
}
void get_new(point &p, double T) {
    p[0] += rnd(rng) * T;
    p[1] += rnd(rng) * T;
}
void codeFire() {
    double T = T0,las = best;
    while(T > 1e-6){
        for(int i = 0; i < L; i++){
            point temp = Point;
            get_new(temp, T);
            double now = calc(temp);
            double DE = now - las;
            if(DE < 0 || exp(-DE / T) > ((rng() % 1000000) / 1000000.0)){
                Point = temp;
                las = now;
                if(now < best){
                    best = now;
                    bestPoint = Point;
                }
            }
        }
        T *= q;
    }
}
int main() {
    std::ios::sync_with_stdio(0),std::cin.tie(0),std::cout.tie(0);
    int n;
    std::cin >> n;
    points.resize(n);
    for (int i = 0; i < n; i++) {
        double x, y, c;
        std::cin >> x >> y >> c;
        points[i] = {x,y,c};
        Point[0] += x,Point[1] += y;
    }
    Point[0] /= n,Point[1] /= n;

    best = calc(Point);
    bestPoint = Point;
    codeFire();
    std::cout << std::fixed << std::setprecision(10) << best << "\n";
    return 0;
}
\end{lstlisting}

\subsection{莫队}
以 $\mathcal O(N \sqrt N)$ 的复杂度完成 $Q$ 次询问的离线查询,
其中每个分块的大小取 $\sqrt N=\sqrt {10^5} = 317$ ,
也可以使用$n / min<int>(n, sqrt(q))$、 $ceil((double)n / (int)sqrt(n))$或者$sqrt(n)$划分。
\begin{lstlisting}[caption={}]
signed main() {
    int n;
    cin >> n;
    vector<int> w(n + 1);
    for (int i = 1; i <= n; i++) {
        cin >> w[i];
    }
    
    int q;
    cin >> q;
    vector<array<int, 3>> query(q + 1);
    for (int i = 1; i <= q; i++) {
        int l, r;
        cin >> l >> r;
        query[i] = {l, r, i};
    }
    
    int Knum = n / min<int>(n, sqrt(q)); // 计算块长
    vector<int> K(n + 1);
    for (int i = 1; i <= n; i++) { // 固定块长
        K[i] = (i - 1) / Knum + 1;
    }
    sort(query.begin() + 1, query.end(), [&](auto x, auto y) {
        if (K[x[0]] != K[y[0]]) return x[0] < y[0];
        if (K[x[0]] & 1) return x[1] < y[1];
        return x[1] > y[1];
    });
    
    int l = 1, r = 0, val = 0;
    vector<int> ans(q + 1);
    for (int i = 1; i <= q; i++) {
        auto [ql, qr, id] = query[i];
        auto add = [&](int x) -> void {};
        auto del = [&](int x) -> void {};
        while (l > ql) add(w[--l]);
        while (r < qr) add(w[++r]);
        while (l < ql) del(w[l++]);
        while (r > qr) del(w[r--]);
        ans[id] = val;
    }
    for (int i = 1; i <= q; i++) {
        cout << ans[i] << endl;
    }
}
\end{lstlisting}

需要注意的是，在普通莫队中，$K$ 数组的作用是根据左边界的值进行排序，当询问次数很少时（$q \ll n$），可以直接合并到 $query$ 数组中。
\begin{lstlisting}[caption={}]
vector<array<int, 4>> query(q);
for (int i = 1; i <= q; i++) {
    int l, r;
    cin >> l >> r;
    query[i] = {l, r, i, (l - 1) / Knum + 1}; // 合并
}
sort(query.begin() + 1, query.end(), [&](auto x, auto y) {
    if (x[3] != y[3]) return x[3] < y[3];
    if (x[3] & 1) return x[1] < y[1];
    return x[1] > y[1];
});
\end{lstlisting}


\subsection{带修莫队}
以 $\mathcal O(N^\frac{5}{3})$ 的复杂度完成 $Q$ 次询问的离线查询，
其中每个分块的大小取 $N^\frac{2}{3}=\sqrt[3]{100000^2}=2154$ （直接取会略快），
也可以使用$pow(n, 0.6666)$划分。
\begin{lstlisting}[caption={}]
signed main() {
    int n, q;
    cin >> n >> q;
    vector<int> w(n + 1);
    for (int i = 1; i <= n; i++) {
        cin >> w[i];
    }
    vector<array<int, 4>> query = {{}}; // {左区间, 右区间, 累计修改次数, 下标}
    vector<array<int, 2>> modify = {{}}; // {修改的值, 修改的元素下标}
    for (int i = 1; i <= q; i++) {
        char op;
        cin >> op;
        if (op == 'Q') {
            int l, r;
            cin >> l >> r;
            query.push_back({l, r, (int)modify.size() - 1, (int)query.size()});
        } else {
            int idx, w;
            cin >> idx >> w;
            modify.push_back({w, idx});
        }
    }
    
    int Knum = 2154; // 计算块长
    vector<int> K(n + 1);
    for (int i = 1; i <= n; i++) { // 固定块长
        K[i] = (i - 1) / Knum + 1;
    }
    sort(query.begin() + 1, query.end(), [&](auto x, auto y) {
        if (K[x[0]] != K[y[0]]) return x[0] < y[0];
        if (K[x[1]] != K[y[1]]) return x[1] < y[1];
        return x[3] < y[3];
    });
    
    int l = 1, r = 0, val = 0;
    int t = 0; // 累计修改次数
    vector<int> ans(query.size());
    for (int i = 1; i < query.size(); i++) {
        auto [ql, qr, qt, id] = query[i];
        auto add = [&](int x) -> void {};
        auto del = [&](int x) -> void {};
        auto time = [&](int x, int l, int r) -> void {};
        while (l > ql) add(w[--l]);
        while (r < qr) add(w[++r]);
        while (l < ql) del(w[l++]);
        while (r > qr) del(w[r--]);
        while (t < qt) time(++t, ql, qr);
        while (t > qt) time(t--, ql, qr);
        ans[id] = val;
    }
    for (int i = 1; i < ans.size(); i++) {
        cout << ans[i] << endl;
    }
}
\end{lstlisting}

\subsection{幻方}
构造题用,其有一些性质(保证$N$为奇数)：$1$ 到 $N^2$
每个数字恰好使用一次,且每行、每列及两条对角线上的数字之和都相同,且为奇数。

构造方式：将 $1$写在第一行的中间,随后不断向右上角位置填下一个数字,直到填满。
\begin{lstlisting}[caption={}]
int n;
cin >> n;
int x = 1, y = (n + 1) / 2;
vector ans(n + 1, vector<int>(n + 1));
for (int i = 1; i <= n * n; i++) {
    ans[x][y] = i;
    if (!ans[(x - 2 + n) % n + 1][y % n + 1]){
        x = (x - 2 + n) % n + 1;
        y = y % n + 1;
    } else {
        x = x % n + 1;
    }
}
\end{lstlisting}

\subsection{约瑟夫问题}
$n$ 个人编号 $0,1,2…,n-1$ ，每次数到 $k$ 出局，求最后剩下的人的编号。

$\mathcal O(N)$ 。
\begin{lstlisting}[caption={}]
int jos(int n,int k){
    int res=0;
    repeat(i,1,n+1)res=(res+k)%i;
    return res; // res+1,如果编号从1开始
}
\end{lstlisting}
$\mathcal O(K\log N)$ ，适用于 $K$ 较小的情况。
\begin{lstlisting}[caption={}]
int jos(int n,int k){
    if(n==1 || k==1)return n-1;
    if(k>n)return (jos(n-1,k)+k)%n; // 线性算法
    int res=jos(n-n/k,k)-n%k;
    if(res<0)res+=n; // mod n
    else res+=res/(k-1); // 还原位置
    return res; // res+1,如果编号从1开始
}
\end{lstlisting}


\subsection{舞蹈链}
理论复杂度大概在$O(c^n)$ 左右，其中 $c$ 为某个非常接近于 $1$的常数,$n$ 为矩阵中 $1$的个数。

数独求解
\begin{lstlisting}[caption={}]
#include<bits/stdc++.h>
constexpr int N = 1e6 + 10;
int ans[10][10], stk[N];

struct DLX {
	static constexpr int MAXSIZE = 1e5 + 10;
	int n, m, tot, first[MAXSIZE + 10], siz[MAXSIZE + 10];
	int L[MAXSIZE + 10], R[MAXSIZE + 10], U[MAXSIZE + 10], D[MAXSIZE + 10];
	int col[MAXSIZE + 10], row[MAXSIZE + 10];

	void build(const int &r, const int &c) {  // 进行build操作
		n = r, m = c;
		for (int i = 0; i <= c; ++i) {
			L[i] = i - 1, R[i] = i + 1;
			U[i] = D[i] = i;
		}
		L[0] = c, R[c] = 0, tot = c;
		memset(first, 0, sizeof(first));
		memset(siz, 0, sizeof(siz));
	}

	void insert(const int &r, const int &c) {  // 进行insert操作
		col[++tot] = c, row[tot] = r, ++siz[c];
		D[tot] = D[c], U[D[c]] = tot, U[tot] = c, D[c] = tot;
		if (!first[r]) first[r] = L[tot] = R[tot] = tot;
		else {
			R[tot] = R[first[r]], L[R[first[r]]] = tot;
			L[tot] = first[r], R[first[r]] = tot;
		}
	}

	void remove(const int &c) {  // 进行remove操作
		int i, j;
		L[R[c]] = L[c], R[L[c]] = R[c];
		for (i = D[c]; i != c; i = D[i])
		for (j = R[i]; j != i; j = R[j])
			U[D[j]] = U[j], D[U[j]] = D[j], --siz[col[j]];
	}

	void recover(const int &c) {  // 进行recover操作
		int i, j;
		for (i = U[c]; i != c; i = U[i])
		for (j = L[i]; j != i; j = L[j]) U[D[j]] = D[U[j]] = j, ++siz[col[j]];
		L[R[c]] = R[L[c]] = c;
	}

	bool dance(int dep) {  // dance
		int i, j, c = R[0];
		if (!R[0]) {
			for (i = 1; i < dep; ++i) {
				int x = (stk[i] - 1) / 9 / 9 + 1;
				int y = (stk[i] - 1) / 9 % 9 + 1;
				int v = (stk[i] - 1) % 9 + 1;
				ans[x][y] = v;
			}
			return true;
		}
		for (i = R[0]; i != 0; i = R[i])
			if (siz[i] < siz[c]) c = i;
		remove(c);
		for (i = D[c]; i != c; i = D[i]) {
			stk[dep] = row[i];
			for (j = R[i]; j != i; j = R[j]) remove(col[j]);
			if (dance(dep + 1)) return true;
			for (j = L[i]; j != i; j = L[j]) recover(col[j]);
		}
		recover(c);
		return false;
	}
} solver;

int GetId(int row, int col, int num) {
	return (row - 1) * 9 * 9 + (col - 1) * 9 + num;
}

void Insert(int row, int col, int num) {
	int dx = (row - 1) / 3 + 1;
	int dy = (col - 1) / 3 + 1;
	int room = (dx - 1) * 3 + dy;
	int id = GetId(row, col, num);
	int f1 = (row - 1) * 9 + num;            // task 1
	int f2 = 81 + (col - 1) * 9 + num;       // task 2
	int f3 = 81 * 2 + (room - 1) * 9 + num;  // task 3
	int f4 = 81 * 3 + (row - 1) * 9 + col;   // task 4
	solver.insert(id, f1);
	solver.insert(id, f2);
	solver.insert(id, f3);
	solver.insert(id, f4);
}
using std::cin;using std::cout;
int main() {
	cin.tie(nullptr)->sync_with_stdio(false);
	solver.build(729, 324);
	for (int i = 1; i <= 9; ++i)
		for (int j = 1; j <= 9; ++j) {
			cin >> ans[i][j];
			for (int v = 1; v <= 9; ++v) {
				if (ans[i][j] && ans[i][j] != v) continue;
				Insert(i, j, v);
			}
		}
	solver.dance(1);
	for (int i = 1; i <= 9; ++i, cout << '\n')
		for (int j = 1; j <= 9; ++j, cout << ' ') cout << ans[i][j];
	return 0;
}
\end{lstlisting}

\subsection{对拍}
\subsubsection{Windows}
\begin{lstlisting}[caption={}]
int main(){
	system("g++ -std=c++20 gen.cpp -o gen.exe");
	system("g++ -std=c++20 brute.cpp -o brute.exe");
	system("g++ -std=c++20 std.cpp -o std.exe");
	for(int i = 1;i <= 1000;i++){
		system("gen.exe > test.in");
		system("brute.exe < test.in > brute.out");
		system("std.exe < test.in > std.out");
		if(system("fc std.out brute.out")){
			cout << "Wrong Answer on test " << i << "\n";
			system("type test.in");
			system("type brute.out");
			system("type std.out");
			return 0;
		}
		else{
			cout << "Accepted on test" << i << ".\n";
		}
	}
	return 0;
}
\end{lstlisting}

\subsubsection{Linux}
\begin{lstlisting}[caption={}]
int main(){
    system("g++ -std=c++20 gen.cpp -o gen");
    system("g++ -std=c++20 brute.cpp -o brute");
    system("g++ -std=c++20 std.cpp -o std");
    for(int i = 1; i <= 1000; i++){
        system("./gen > test.in");
        system("./brute < test.in > brute.out");
        system("./std < test.in > std.out");
        if(system("diff std.out brute.out")){
            cout << "Wrong Answer on test " << i << "\n";
            system("cat test.in");
            system("cat brute.out");
            system("cat std.out");
            return 0;
        }
        else{
            cout << "Accepted on test " << i << ".\n";
        }
    }
    return 0;
}
\end{lstlisting}

\subsection{python高精度}
\begin{lstlisting}[caption={}]
t decimal
decimal.getcontext().prec = 1000000
n = decimal.Decimal(input())
t = (4 + n) * (3 + n) * (2 + n) * (1 + n)
print(t // 24)
\end{lstlisting}


\subsection{乱七八糟}
看数据范围（多测总和）,变量shadow,多测清空,
longlong,数组开小了,模数不对,MLE?,对拍记得看输出在不在变,
输出格式,inf开小,答案初值,STL重构导致引用失效,极端情况(n=1),
删除调试代码,循环内变量名重复\\

将二维矩阵旋转90°: $b[j][n - j + 1] = a[i][j]$\\
将二维矩阵旋转180°: $b[n - i + 1][n - j + 1] = a[i][j]$\\
将二维矩阵旋转270°: $b[n - j + 1][i] = a[i][j]$\\
